\section{Methodology}

\subsection{Bayesian Heirarchical Models}

A hierarchical model is a statistical model that includes multiple levels of parameters, where the parameters at one level are treated as random variables that depend on parameters at a higher level. In a Bayesian hierarchical model, we specify prior distributions for the parameters at each level, and use Bayes' theorem to update our beliefs about these parameters based on observed data.

We would like to capture both individual battery characteristics and the influence of different usage (cycling) conditions. A two level hierarchical structure is considered:
\begin{itemize}
    \item Level 1 (Cell-level parameters $\theta_j$): Each battery $j$ has its own parameters $\theta_j$ that represent how early cycle features relate to battery aging for that particular subgroup or cycling condition.
    \item Level 2 (Hyperparameters $\gamma$): These represent the parameters that govern the distribution of cell-level parameters across all subgroups, meaning they capture similarities across different batteries and cycling conditions.
\end{itemize}

Bayes rule is used to combine prior knowledge with the observed early-cycle data from multiple batteries to estimate posterior distributions of parameters at both levels. 

The level 2 posterior is calculated as 
$$
P(\gamma | \{Y\}) \propto \text{Likelihood} \times \text{Prior}
$$

\subsection{Multichain Monte Carlo}